% ---------------------------------------------------------------------------
% MSC PROPOSAL V1.4 — Cost-Aware Repository Change Localization with Sparse
% Impact Planning and Bounded Verification
%
% Historical version (V1.4). V1, V1.2 and V1.3 are preserved immutably.
% V1.4 integrates, conservatively: (1) the completed Saleor DEVELOPMENT
% Sparse run (150x3, 446 valid / 4 failed, 7,316,986 tokens / $2.31,
% 0 truncations); (2) the Saleor frozen Route-B transfer replication
% (149 evaluable tasks; B=5 composite 0.237 vs analytic Random 0.006,
% delta +0.231 CI [+0.180,+0.287]; 5/5 folds positive; DEVELOPMENT transfer
% replication, NOT confirmatory); (3) the ranker-identity audit
% (the frozen arm historically labelled "Classical-CIA" is exactly
% normalized BM25 + binary graph-neighbor; the Hybrid arm is a redundant
% rank-equivalent alias); (4) the incremental-evidence ablation showing the
% cross-repo signal is predominantly lexical (BM25) with a small,
% largely non-significant graph contribution.
% ---------------------------------------------------------------------------
\documentclass[11pt,a4paper]{article}
\usepackage[T1]{fontenc}
\usepackage[utf8]{inputenc}
\usepackage{lmodern}
\usepackage{microtype}
\usepackage{amsmath,amssymb}
\usepackage{booktabs,tabularx,array}
\usepackage{enumitem}
\usepackage{graphicx}
\usepackage[margin=2.5cm]{geometry}
\usepackage{xcolor}
\usepackage{url}
\usepackage[colorlinks=true,linkcolor=blue,citecolor=blue,urlcolor=blue]{hyperref}
\emergencystretch=1.5em

\title{\textbf{Cost-Aware Repository Change Localization with Sparse Impact\\[2pt]
Planning and Bounded Verification}\\[6pt]
\large MSc Research Proposal (V1.4, supervisor-ready / print-candidate)}
\author{Ahmed Ehab}
\date{2026-09-17}

\begin{document}
\maketitle

\begin{abstract}
Large language models (LLMs) that edit a software repository must first decide
which files a requested change affects. This decision is costly: an explicit
file-by-file impact policy over a repository-scale candidate universe spends
hundreds of serialized records per call, and a second-stage reasoner that
revisits the scope can multiply inference cost by orders of magnitude.

\emph{Central identity of the thesis:}

\begin{quote}
This thesis studies whether a cheap sparse file-impact policy can reduce
repository-wide localization effort while a bounded second look recovers
important files the first pass chose not to change.
\end{quote}

We investigate a representation that is cheap by construction ---
\emph{preserve-by-omission} sparse impact plans --- and evaluate it with cheap
lexical/static baselines and an explicit \emph{bounded} second-stage
omission-recovery mechanism, all under matched inference budgets. Preliminary
development evidence shows a large serialization-cost reduction for the sparse
representation, a meaningful zero-LLM lexical signal, and --- for the first
time at candidate level --- a robust recovery signal: a frozen
BM25+Graph-Neighbor composite ranker (historical label: Classical-CIA)
recovers Sparse-omitted missed files above the analytic-Random baseline across
the whole budget curve, with task-level bootstrap intervals excluding zero.
A clean second-repository DEVELOPMENT run (Saleor, 150 tasks) transfers
positively: the frozen composite beats analytic Random on 149 evaluable tasks
(B=5 delta +0.231, CI [+0.180,+0.287]), a DEVELOPMENT transfer replication ---
not a confirmatory result. The confirmatory run on the authorized djangoCMS
INTERNAL\_TEST (80 tasks, 2026-09-17) **CONFIRMS** the fixed-budget recovery
signal at every B (B=5 composite 0.165 vs analytic Random 0.028, delta +0.137,
CI [+0.075,+0.205]). An identity audit and incremental-evidence ablation
show the signal is predominantly lexical (BM25), with the binary graph-neighbor
term contributing a small, largely non-significant increment; the thesis
therefore does not claim graph novelty. Task-level omission risk is not
reliably predictable on current development evidence, so the thesis does not
depend on a task-level risk scorer.
\end{abstract}

% ---------------------------------------------------------------------------
\section{Background, motivation, and problem}
\label{sec:background}

Repository-level coding agents and impact planners must select affected files
before editing. Two problems motivate the thesis. First, \emph{impact
inference}: did the model identify the right files? Second,
\emph{impact-policy representation}: can a complete file-action policy be
expressed without spending output budget on hundreds of repeated
\emph{preserve} decisions? Our controlled study shows that an explicit full
policy at a 16K completion cap serializes on average 144 decision records per
candidate universe, while a sparse policy serializes on average 4.9 records (a
$\sim$96.6\% reduction), with large associated reductions in completion tokens
and cost.

Second-stage reasoning is expensive: a graph-guided agent evaluated under a
shared protocol consumed on the order of a million tokens per task on the same
held-out tasks where our sparse first pass stays below a few thousand tokens.
If a cheap first pass already localizes the bulk of affected files, the
question is whether a \emph{bounded} second look can recover missed files
without paying the full always-on cost.

Given a repository revision $P$ (parent commit), a change request $q$, and a
candidate universe $U(P)$ visible at $P$, a localization system produces an
affected-file set $\hat{S}(P,q)$; the hidden observed change set $S^*(P)$ is
used for evaluation only. The observed historical diff is an \emph{observed
change-set proxy}, never semantic impact gold (Section~\ref{sec:dataset}).

% ---------------------------------------------------------------------------
\section{Central thesis hypothesis and gap}
\label{sec:hypothesis}

\emph{Hypothesis.} A sparse/cheap first-pass localization policy followed by a
bounded omission-recovery mechanism can dominate both cheap-localization-only
and always-on expensive reasoning on the correctness--cost Pareto frontier, for
the repositories and protocols we evaluate.

Everything else supports this story:
\begin{itemize}
  \item \emph{Preserve-by-Omission} = representation / first-pass contribution;
  \item \emph{BM25 / classical static change-impact analysis} = cheap baselines;
  \item \emph{graph-guided agent} = expensive system-level context;
  \item \emph{task-level omission-risk scorer} = investigated and currently
        rejected on development evidence;
  \item \emph{candidate-level bounded verification} = primary next mechanism;
  \item \emph{second and third repositories} = generalization.
\end{itemize}

Task-level RiskScorer success is \emph{not} a thesis requirement. Existing work
covers several individual components; this thesis evaluates \emph{their
combination} as a sparse first pass plus bounded omission recovery under
matched budgets (Section~\ref{sec:related}).

% ---------------------------------------------------------------------------
\section{Research questions and expected contributions}
\label{sec:rq}

\begin{description}
  \item[RQ1.] How does sparse explicit impact-policy representation affect
  inference/output cost and file-selection fidelity?
  \item[RQ2.] How accurately can cheap lexical/static methods localize impacted
  files relative to LLM planners under shared protocols?
  \item[RQ3.] Can bounded second-stage verification recover missed impacted
  files more efficiently than always-on or random matched-budget verification?
  \item[RQ4.] How do these correctness--cost trade-offs transfer across
  repositories and ecosystems?
\end{description}

If task-level risk modeling remains unsupported, RQ3 is answered by the
candidate-level mechanism. \emph{Expected contributions:}
(1) sparse explicit impact-plan representation with controlled evidence;
(2) a reproducible real-commit impact-localization dataset and protocol;
(3) fair cheap/classical/LLM/agent comparison under common file-level metrics
and explicit budgets;
(4) a bounded/selective-compute verification protocol (candidate-level by
default);
(5) cross-repository correctness--cost evidence;
(6) reproducibility artifacts and leakage-control methodology.

No positive algorithmic result is promised.

% ---------------------------------------------------------------------------
\section{Positioning against related work}
\label{sec:related}

Existing work covers several individual components. This thesis evaluates their
combination. Verified related lines (primary sources, 2026-09-16/17):

\begin{itemize}
  \item \emph{Classical change-impact / ripple effect (RIPPLE line).}
  RIPPLE = a known seed $\rightarrow$ recall expansion $\rightarrow$ precision
  refinement \cite{haney1972_ripple}. This thesis differs: an explicit
  repository-wide sparse policy $\rightarrow$ a rejected residual set
  $\rightarrow$ bounded reconsideration of rejected candidates.
  \item \emph{Change-Patterns Mapping} (TSE 2022) \cite{huang2022changepatterns}:
  boosts/reranks the initial impact set from historical patterns. This thesis
  allocates \emph{new inspection to the excluded residual set} rather than
  reranking the selected set.
  \item \emph{ArtifactSync}: a demo-level mechanistic analogue (uncertainty-
  driven deeper context) but a different post-commit synchronization task, not
  file-impact localization.
  \item \emph{Modern transform + program-dependence-graph impact analysis}
  \cite{yan2026transformerpdg}: graph+LLM impact analysis is active; our graph
  is one optional verifier.
  \item \emph{Selective-compute / efficiency+verification priors}
  \cite{wu2024repoformer,zhao2025fastcoder}: selective retrieval / efficient
  verification at repository level for code-completion/generation tasks,
  guarding generic ``selective escalation is novel'' claims.
  \item \emph{Repository-level coding and retrieval agents}
  \cite{agentless2024,codeplan2024,repocoder2023,autocoderover2024,repoagent2024}:
  LLM agents and retrieval-conditioned generation over whole repositories;
  this thesis evaluates a cheap sparse first pass against shared-cost context.
  \item \emph{Graph-guided issue localization}
  \cite{locagent2025,graphlocator2026,jiang2025issueloc}: expensive graph-guided
  competitors (LocAgent, GraphLocator, LLM-driven iterative code-graph search).
  \item \emph{Classical / history-learned change-impact predictors}
  \cite{abdeen2015dependency,iwasaki2022impacttool,borg2017recommendationcia,etezadi2025llmcia}:
  dependency- and history-based impact prediction and industrial studies.
\end{itemize}

Graph, history, and selective computation are \emph{not} novelty claims by
themselves. The comparison matrix (Section~\ref{sec:design}) records whether
each line has an explicit file policy, sparse omission representation, cheap
first pass, omission recovery, hard/matched budget, real-commit evaluation, and
cross-repository evidence. Where novelty is only defensible once experiments
support it, it is marked \emph{candidate novelty --- not yet claimed}.

% ---------------------------------------------------------------------------
\section{Proposed two-stage architecture}
\label{sec:architecture}

\begin{enumerate}
  \item \textbf{Sparse/cheap first-pass impact policy}: explicit file-level
        policy with preserve-by-omission serialization.
  \item \textbf{Candidate-level bounded verification} (primary): rank the
        first-pass omitted candidates with independent structural / retrieval /
        history evidence, and verify only those under a hard per-task budget
        $B$. Compare with always-verify and random matched-budget verification.
\end{enumerate}

\begin{center}
\begin{tabular}{@{}c@{}}
\hline
Change request \\
$\downarrow$ \\
Sparse / cheap first pass \\
$\downarrow$ \\
Omitted candidates \\
$\downarrow$ \\
Suspicion / ranking mechanism \\
$\downarrow$ \\
Bounded verification (budget $B$) \\
$\downarrow$ \\
Final affected-file set \\
\hline
\end{tabular}
\end{center}

% ---------------------------------------------------------------------------
\section{Dataset construction, labels, and leakage controls}
\label{sec:dataset}

\begin{itemize}
  \item \textbf{Real history, parent-only inputs.} Inference uses the parent
        commit $P$; the target $T$ is hidden. Live HEAD is never used.
  \item \textbf{Eligibility and exclusions.} Single-parent commits with a
        meaningful intent, touching only production source files; tests,
        configurations, docs, migrations, and generated files are excluded;
        intent-path leakage is a hard exclusion.
  \item \textbf{Deduplication.} Exact-proxy-set / shared-reference / suspected-
        related changes are removed with frozen deterministic rules.
  \item \textbf{Observed change-set proxy.} $P\!\to\!T$ changed production
        files are an \emph{observed} proxy, never semantic impact gold.
        Localization claims are scoped to historical-file recovery unless
        semantic adjudication supports more.
  \item \textbf{Splits.} Metadata-only, deterministic split frozen before any
        model result. For djangoCMS V2-LARGE: development, an untouched
        internal test, and a reserve. The v1 40 are LEGACY-EXPOSED and never
        reused as untouched confirmation.
\end{itemize}

% ---------------------------------------------------------------------------
\section{Experimental design, baselines, budgets, metrics, statistics}
\label{sec:design}

\subsection{Budget curve, not a single point}
The primary scientific object is the \emph{budget curve} $B \in \{0,1,3,5,10\}$;
$B{=}0$ is Sparse alone. $B{=}5$ is a convenient reference operating point, not
a scientifically privileged single claim. We also report the normalized
sensitivity $B / |\text{omitted set}|$ as a secondary view.

\subsection{Analytic Random control}
For ranking-only experiments the Random control is the analytic hypergeometric
expectation: for a task with $N_t$ omitted candidates, $M_t$ Sparse-observed
FNs, and budget $B$, $\mathbb{E}[X_t] = B \cdot M_t / N_t$ (with $B$ clipped to
$N_t$). The comparison is per-task observed recovery vs this analytic
expectation, aggregated with task-level bootstrap intervals. Random repetitions
are never independent tasks.

\subsection{Primary endpoint}
\emph{False-Negative Recovery Rate / Omission Recovery Rate} @ $B$: per task,
recovered Sparse-observed FNs within $B$ / all Sparse-observed FNs. Secondary:
absolute recall gain, FNR reduction, fraction of Oracle@$B$ gap closed,
candidates inspected per recovered FN, precision/F1 (safety metrics), runtime,
index/build cost.

\subsection{Anchors (per B)}
Sparse/$B{=}0$; Analytic Random; BM25; Path-token; Graph/static;
BM25+Graph-Neighbor Composite (historical label: Classical-CIA; the frozen
primary); history/co-change (where parent-visible history is available); one
fixed simple hybrid (0.5 BM25 + 0.5 graph neighbor, a rank-equivalent
redundant control --- confirmed by the ranker-identity audit, not an
independent baseline); Oracle (ranking headroom); InspectAll (exhaustive).
\emph{Oracle@B $\neq$ InspectAll}: Oracle is ranking headroom;
InspectAll is exhaustive reconsideration.

\subsection{Experimental Design at a Glance}
\begin{table}[h]
\centering
\small
\renewcommand{\arraystretch}{1.3}
\setlength{\tabcolsep}{3pt}
\begin{tabularx}{\textwidth}{@{}>{\raggedright\arraybackslash}p{1.7cm}>{\raggedright\arraybackslash}p{1.8cm}>{\raggedright\arraybackslash}p{1.8cm}>{\raggedright\arraybackslash}X>{\raggedright\arraybackslash}p{1.1cm}>{\raggedright\arraybackslash}p{1.6cm}>{\raggedright\arraybackslash}X>{\raggedright\arraybackslash}X@{}}
\toprule
Study & Purpose & Repo / Split & Method & Unit & Budget / Reps & Primary metrics & Status \\
\midrule
Controlled & sparse-vs-full cost & djangoCMS / controlled & Full vs Sparse & task & cap 16K / 3 reps & tokens, records, cost, P/R/F1 & Done \\
Historical & selection fidelity & djangoCMS / held-out & Full vs Sparse & task & cap 16K / 3 reps & P/R/F1/FNR & Done \\
Agent & expensive context & djangoCMS / held-out & graph-guided agent & task & per-call & P/R/F1, cost & Done \\
Cheap baselines & zero-LLM signal & djangoCMS / dev & BM25 / graph & task & K curve & P/R/F1 & Done \\
Task risk & risk predictability & djangoCMS / dev & features & task & 3 reps & AUROC\slash AUPRC & Negative / closed \\
Candidate recovery & omission recovery & djangoCMS / dev & composite ranker + verifier & cand. & $B\in\{0..10\}$ & FN-recovery & DEV replicated; confirmatory pending \\
Cross-repo transfer & transfer & Saleor / dev & frozen method & task & B curve & per-repo & Saleor DEV replication complete \\
\bottomrule
\end{tabularx}
\caption{Experimental Design at a Glance (internal IDs appear in traceability notes only). Status as of 2026-09-17.}
\end{table}

\subsection{Statistics and multiple-comparison discipline}
Task is the independent unit; candidate rows are nested. One primary endpoint
(recovery@$B$); one primary arm-vs-Random comparison; secondary endpoints
labeled; task-level paired bootstrap; no post-hoc feature-family winner;
exploratory analyses reported as exploratory.

% ---------------------------------------------------------------------------
\section{Preliminary evidence}
\label{sec:prelim}

All numbers are from audited studies in this repository; they are development
or controlled evidence with their bounds.

\begin{itemize}
  \item \textbf{Representation effect (controlled).} At a 16K completion cap,
        Sparse serialized on average \textbf{4.9} decision records vs
        \textbf{144.0} for Full ($\sim$\textbf{96.6\%} reduction), with large
        completion-token and cost reductions. No universal semantic-superiority
        claim.
  \item \textbf{Historical real-commit study.} 60 cells over 10 held-out real
        djangoCMS commits: selection is difficult, FNR is high, and there is no
        Full-vs-Sparse superiority claim (paired deltas' intervals cross zero).
        Sparse is substantially cheaper.
  \item \textbf{Graph-guided agent comparison.} System-level shared-protocol
        context: 402 calls / 32.8M tokens / $\sim$\$9.93 estimated on 10 tasks,
        with 5/10 non-usable outcomes (2 timeout + 1 context-length + 2
        completed-but-empty). Not a reproduction of the published fine-tuned
        result.
  \item \textbf{Cheap/static development baselines.} BM25 provides a meaningful
        zero-LLM signal; graph-ranked seeds are comparable; a classical CIA
        (dependency propagation) baseline is implemented on 150 development
        cases.
  \item \textbf{Task-level omission-risk investigation.} Live Sparse
        development inference (90 cells; then 431 calls with a fail-closed stop
        at the 2.5M-token ceiling). A task-level RiskScorer is not justified:
        the v1 signal does not replicate within the new development cohort and
        the pooled signal is a candidate-universe-size artifact.
  \item \textbf{Candidate-level bounded verification (Route B, development).}
        A frozen BM25+Graph-Neighbor composite ranker (historical label:
        Classical-CIA; exact formula = normalized BM25 + binary graph-neighbor,
        per the ranker-identity audit) recovers Sparse-omitted missed files
        above the analytic-Random baseline across the whole budget curve
        $B \in \{1,3,5,10\}$, with task-level bootstrap intervals excluding
        zero at every point (e.g.\ at $B{=}5$: pooled macro recovery 0.163 vs
        random 0.029, interval $[+0.09,+0.18]$), 5/5 development folds
        positive, and no omitted/universe-size artifact. A small verifier
        pilot (30 calls, $\sim$\$0.002) found the frozen ranker places the
        recoverable missed files within the top-$B$ (Oracle-in-top-$B$
        $\approx$ 1.0) and the verifier recovers $\sim$0.86--1.0 of them, so
        the dominant loss is first-pass omission.
  \item \textbf{Saleor DEVELOPMENT transfer replication (2026-09-17).} A clean
        Saleor DEVELOPMENT Sparse run (150 tasks $\times$ 3 reps = 450 cells;
        446 valid / 4 failed; 7,316,986 tokens / $\sim$\$2.31; 0 truncations;
        no Saleor-specific tuning) enabled a frozen Route-B transfer test on
        149 evaluable tasks: the composite beats analytic Random at every
        $B$ ($B{=}5$: 0.237 vs 0.006, delta +0.231, CI
        $[+0.180,+0.287]$; 5/5 folds positive; no size artifact). This is a
        \textbf{DEVELOPMENT transfer replication, NOT a confirmatory result};
        Saleor INTERNAL\_TEST/RESERVE remain sealed.
  \item \textbf{djangoCMS confirmatory result (2026-09-17, CONFIRMS).} The
        frozen confirmatory protocol ran on the authorized djangoCMS V2
        INTERNAL\_TEST (80 tasks) under the frozen V2 packet and budget
        (560 calls / 1,470,174 tokens / $\sim$\$0.51; 0 failures; 0 excluded):
        the composite beats analytic Random at every B ($B{=}1$ 0.059 vs
        0.006; $B{=}3$ 0.110 vs 0.017; $B{=}5$ 0.165 vs 0.028, delta +0.137,
        CI $[+0.075,+0.205]$; $B{=}10$ 0.267 vs 0.055), with task-level
        bootstrap intervals excluding zero at every B and no size artifact.
        **Fixed-budget Route B is scientifically CONFIRMED on djangoCMS
        INTERNAL\_TEST.** The confirmatory test is permanently used and will
        not be reused for adaptive-budget (P2) selection.
  \item \textbf{Incremental-evidence ablation (zero API, 2026-09-17).}
        Composite-minus-BM25 paired task-level deltas are small with bootstrap
        CIs including zero at most $B$ on both repositories (djangoCMS
        $B{=}5$ +0.003 $[-0.015,+0.019]$; Saleor $B{=}5$ +0.003
        $[-0.017,+0.024]$); the replicated cross-repo signal is
        \textbf{predominantly lexical (BM25)}. The thesis does not claim graph
        novelty.
  \item \textbf{Classical/history evidence.} A parent-visible history/co-change
        arm also beats analytic Random where history is available; the
        BM25+Graph-Neighbor composite remains the frozen primary.
\end{itemize}

These results motivate a larger development dataset, protection of an untouched
internal test, and the bounded-verification design; they are not rewritten as a
success.

% ---------------------------------------------------------------------------
\section{Cross-repository generalization}
\label{sec:gen}

Stage 1 djangoCMS V2-LARGE (within-repo large evidence + untouched internal
test) $\to$ Stage 2 Saleor (different real system, same Python/Django stack)
$\to$ Stage 3 NestJS (cross-language / cross-framework evidence). Each stage
states its claim boundary. NestJS is preferred but not forced; a predeclared
eligible-pool yield criterion and a backup protocol are frozen.

\subsection{Saleor DEVELOPMENT transfer replication (completed, 2026-09-17)}
The Saleor DEVELOPMENT sparse run (150$\times$3 = 450 cells; 446 valid / 4
failed; 7,316,986 tokens / $\sim$\$2.31; 0 truncations; no Saleor-specific
tuning) and the frozen Route-B transfer test (149 evaluable tasks) replicate
the bounded omission-recovery signal on a second, much larger repository
($B{=}5$ composite 0.237 vs analytic Random 0.006; delta +0.231, CI
$[+0.180,+0.287]$; 5/5 folds positive; no size artifact). This is
\textbf{DEVELOPMENT transfer replication only --- NOT confirmatory}. Saleor
INTERNAL\_TEST/RESERVE are sealed and were never opened. Per-repository primary;
no djangoCMS+Saleor pooling as the headline.

\subsection{djangoCMS confirmatory result (2026-09-17, CONFIRMS)}
The frozen confirmatory protocol ran on the authorized djangoCMS V2
INTERNAL\_TEST (80 tasks) under the frozen V2 packet and budget (560 calls /
1,470,174 tokens / $\sim$\$0.51; 0 failures; 0 excluded): the frozen
BM25+Graph-Neighbor Composite recovers Sparse-omitted historically-changed
files above the analytic Random control across the whole budget curve
($B{=}1$ 0.059 vs 0.006; $B{=}3$ 0.110 vs 0.017; $B{=}5$ 0.165 vs 0.028;
$B{=}10$ 0.267 vs 0.055), with task-level bootstrap intervals excluding zero
at every B and no size artifact. **The fixed-budget Route-B result is
scientifically CONFIRMED on djangoCMS INTERNAL\_TEST.** The confirmatory
test is permanently used and will not be reused for adaptive-budget (P2)
selection; P2 is evaluated on development data only, with Saleor
INTERNAL\_TEST sealed as a possible future P2 confirmation set.

\subsection{Scope contract}
\emph{CORE MSc:} Preserve-by-Omission representation evidence; fixed-budget
Route-B development + the confirmed djangoCMS INTERNAL\_TEST result;
semantic-proxy audit; djangoCMS; Saleor as secondary replication if completed.
\emph{Conditional / Future:} adaptive $B_t$ (P2, development-data exploration,
Nov 2026 -- Mar 2027); learned budget allocation; NestJS cross-language;
end-to-end regeneration correctness; persistent memory; richer causal-graph
verifier.

% ---------------------------------------------------------------------------
\section{Threats, contingency, and timeline}
\label{sec:threats}

\subsection{Threats to validity}
Internal: the observed proxy is not semantic gold; leakage control is mandatory;
no hidden-test peeking. External: $n$ and single-repository generalization; V2
and cross-repository evidence are the mitigation, with per-repository primary
metrics. Statistical: low negatives, multiple comparisons, wide random bands ---
addressed by pre-registration, class-balance gates, analytic-Random controls,
and bootstrap uncertainty. Representation: LLM serialization cost spans
providers; both recorded and frozen-pricing cost are reported.

\subsection{Semantic-proxy validity}
Because the historical diff is only a proxy, we prepare a stratified
development-only sample with evidence packets for later human adjudication.
\begin{sloppypar}
The protocol is
\texttt{docs/SEMANTIC\_PROXY\_AUDIT\_PROTOCOL.md}.
\end{sloppypar}
Machine-assisted notes are never semantic gold.

\subsection{Contingency}
If candidate-level bounded verification does not beat analytic Random at the
frozen budget, the negative result is frozen and reported; no complex verifier
is invented.

\subsection{Future work: adaptive / cost-sensitive verification allocation}
Inspired by the Shichao Zhang adaptive-resource line (Learning-$k$, one-step
KNN, cost-sensitive KNN, adaptive-neighborhood)
\cite{zhang_tkde2022_challenges,zhang_tkde2021_onestep,zhang_neurom2020_costsensitive,zhang_tkde2022_reachable}
and selective-prediction / learning-to-defer literature, we sketch a future
adaptive budget: $B_t^* = \arg\min_B [\,C_{\text{verify}}(B) + \rho\,\mathbb{E}[\text{remaining omission
loss}\mid \text{evidence}, B]\,]$, with $\rho$ a sensitivity, not a claimed
monetary truth. This is NOT yet a contribution and is NOT implemented before
fixed-budget Route B is established (status: CONDITIONAL; see the P2
pre-registration note).

\subsection{Timeline (2026-11 to 2027-10)}
\begin{table}[h]
\centering
\small
\renewcommand{\arraystretch}{1.25}
\setlength{\tabcolsep}{6pt}
\begin{tabularx}{\textwidth}{@{}>{\raggedright\arraybackslash}p{1.9cm}>{\raggedright\arraybackslash}X@{}}
\toprule
\textbf{Window} & \textbf{Planned work} \\
\midrule
2026-11 & Fixed Route-B confirmation (djangoCMS INTERNAL\_TEST; CONFIRMED 2026-09-17) + P2 literature/formulation freeze \\
2026-12 & Simple Shichao-inspired adaptive-budget policies (score-gap, marginal-score, cost-ratio stopping) on development data \\
2027-01 & Cost-sensitive / one-step / demand-driven adaptive-budget analogues (development data only) \\
2027-02 & Newly discovered competitor/alternative algorithm reproduction (selective/cascade/optimal-stopping families) \\
2027-03 & P2 comparison/freeze/pre-registration on development; negative closure if not justified; Saleor INTERNAL\_TEST only if pre-registered \\
2027-04 & Optional external-validity extension (NestJS or predeclared backup if suitability gate passes); no forced cross-language claim \\
2027-05 & Consolidated statistical analysis; Pareto/frontier analysis; semantic-audit integration; final threat-to-validity revision \\
2027-06 & Thesis chapters/results writing; artifact cleanup; reproducibility package; supervisor review round 1 \\
2027-07 & Thesis full draft; corrections; final experimental closure; submission-ready research manuscript draft \\
2027-08 & Target research/thesis completion; supervisor review round 2; final figures/tables/appendices; defense/seminar preparation \\
2027-09 & Publication revision/submission window; reviewer-style internal audit; contingency for delayed replication or manuscript resubmission \\
2027-10 & Final buffer: journal/conference resubmission, thesis formatting/administrative closure, archival artifact release \\
\bottomrule
\end{tabularx}
\caption{Work schedule, November 2026 -- October 2027. July--August 2027 is the intended substantive research/thesis completion; September--October 2027 is an honest publication/revision/administrative buffer, not fabricated work. November 2026 -- March 2027 is the P2 adaptive-budget exploration program (development data only).}
\end{table}
Intended substantive research/thesis completion: \textbf{2027-07/08}. The
2027-09/10 window is a publication, revision, and administrative buffer. No
experiments are fabricated merely to fill months. P2 (adaptive budget) is a
conditional exploration: if it does not justify a candidate on development
data, it is closed as negative and the fixed Route-B result (already confirmed
on djangoCMS INTERNAL\_TEST) remains the thesis basis. NestJS is conditional on
its frozen gates and evidence.

% ---------------------------------------------------------------------------
\section{References}
\label{sec:refs}

All entries below were verified against primary sources (arXiv API /
Crossref) on 2026-09-16/17. Preprints are explicitly marked. No fabricated or
placeholder citations remain; the source bibliography is
\texttt{references.bib} (V1.4).

\bibliographystyle{unsrt}
\bibliography{references}

\end{document}